\chapter{Building Agents: Strands and LangGraph}

The previous chapters established the conceptual architecture of an agent: layers of
intelligence, action, memory, and security. This chapter makes those abstractions concrete
by implementing the same assistant twice, once with Strands and once with LangGraph. The
tools, the gateway, the skills, and the memory resource are identical in both cases. Only
the way the agent loop is assembled differs. Examining both implementations in parallel
surfaces what each framework provides and what it leaves to the developer.

The assistant connects to backend capabilities through an MCP gateway secured with AWS
Signature Version 4 (SigV4). A \texttt{SKILL.md} document encodes the operating procedure.
AgentCore Memory keeps track of conversational state and user preferences across sessions.
The book follows the order in which components are typically written: tools first, then the
agent loop, then memory, then the entrypoint.

\section{Step 1: Define the tools}

Tools are the agent's interface to the world. Before picking a framework, define what the
agent can do. Built-in tools are plain Python functions decorated with the \texttt{@tool}
decorator. The Strands import and the LangGraph import use different packages, but the
pattern is identical: a function, a docstring, and an optional return type.

\begin{lstlisting}[language=Python]
# Strands import
from strands import tool

# LangGraph import
from langchain_core.tools import tool

@tool
def get_current_time() -> dict[str, str]:
    """Return the current UTC time."""
    from datetime import datetime, timezone
    return {"utc": datetime.now(timezone.utc).isoformat()}

@tool
def echo(message: str) -> dict[str, str]:
    """Return the message back to the caller. Useful for testing."""
    return {"response": message}
\end{lstlisting}

These two tools illustrate the pattern: a decorated function, a docstring that the model
uses to decide when to call the tool, and a typed return value. In practice, built-in tools
handle lightweight operations that do not require a gateway round-trip---simple lookups,
formatting, or diagnostic checks. The more capable tools are loaded dynamically from the
AgentCore Gateway at startup, as described in the next section.

\section{Step 2: Authenticate to the gateway with SigV4}

Because the AgentCore Gateway enforces IAM authentication on every inbound request, the
agent must sign each HTTP call with AWS Signature Version 4 before sending it. This is the
standard authentication pattern for any AWS workload calling an IAM-protected service
endpoint---the same pattern used when a container calls S3, DynamoDB, or any other AWS API
directly. Since the MCP client uses \texttt{httpx} for transport, the signing logic is
implemented as a custom \texttt{httpx.Auth} subclass that resolves the container's
execution role credentials via \texttt{botocore} and attaches the required headers to each
request.

It is important to note that this signing layer covers only the agent-to-gateway leg of the
call. Once the request reaches the gateway, authentication to backend tools such as Lambda
functions is handled by the gateway itself using its own service role. The agent code is
never responsible for signing those outbound calls.

\begin{lstlisting}[language=Python]
import hashlib
from urllib.parse import urlparse

import botocore.session
import httpx
from botocore.auth import SigV4Auth
from botocore.awsrequest import AWSRequest

SERVICE_NAME = "bedrock-agentcore"

class SigV4HTTPAuth(httpx.Auth):
    requires_request_body = True

    def __init__(self, service: str, region: str,
                 credentials=None):
        self.service = service
        self.region = region
        self.credentials = credentials

    def auth_flow(self, request: httpx.Request):
        creds = self._resolve_creds()
        body = bytes(request.content) if request.content else b""
        signed_headers = self._sign(request, body, creds)
        yield httpx.Request(
            method=request.method,
            url=request.url,
            headers=signed_headers,
            content=body,
        )

    def _resolve_creds(self):
        if self.credentials:
            return self.credentials
        session = botocore.session.get_session()
        creds = session.get_credentials()
        if not creds:
            raise RuntimeError("AWS credentials not found")
        return creds

    def _sign(self, request, body, creds):
        parsed = urlparse(str(request.url))
        headers = {
            "Host": parsed.netloc,
            "Content-Type": "application/json",
            "Content-Length": str(len(body)),
            "X-Amz-Content-Sha256": hashlib.sha256(
                body
            ).hexdigest(),
        }
        aws_req = AWSRequest(
            method=request.method,
            url=str(request.url),
            data=body,
            headers=headers,
        )
        SigV4Auth(
            creds, self.service, self.region
        ).add_auth(aws_req)
        return dict(aws_req.headers)
\end{lstlisting}

This class is shared across both framework implementations. Any HTTP request sent through
this auth handler has its body hashed, its headers signed, and its timestamp included
before it leaves the agent container.

\section{Step 3: Load gateway tools over MCP}

Gateway tools arrive via the MCP protocol. The two frameworks use different MCP clients,
so the gateway loading code diverges here. Strands uses \texttt{MCPClient} from the
\texttt{strands.tools.mcp} package and loads tools synchronously during initialisation.
LangGraph uses \texttt{MultiServerMCPClient} from \texttt{langchain-mcp-adapters} and
loads tools asynchronously.

\paragraph{Strands gateway loader.}

\begin{lstlisting}[language=Python]
import os
import atexit
from mcp.client.streamable_http import streamablehttp_client
from strands.tools.mcp import MCPClient

def _gateway_urls() -> list[str]:
    raw = os.environ.get("GATEWAY_URL", "")
    return [url.strip() for url in raw.split(",") if url.strip()]

def _make_client(url: str) -> MCPClient:
    region = os.environ.get("GATEWAY_REGION", "us-west-2")
    auth = SigV4HTTPAuth(service=SERVICE_NAME, region=region)
    return MCPClient(
        lambda: streamablehttp_client(
            url=url,
            headers={
                "Accept": "application/json, text/event-stream"
            },
            auth=auth,
        )
    )

class GatewayToolPool:
    def __init__(self) -> None:
        self.clients: list[MCPClient] = []
        self.tools_list: list = []
        self.load_tools()
        atexit.register(self.stop)

    def load_tools(self) -> None:
        urls = _gateway_urls()
        clients = []
        for url in urls:
            client = _make_client(url)
            client.__enter__()
            clients.append(client)
        tools = []
        for client in clients:
            tools.extend(client.list_tools_sync())
        self.clients = clients
        self.tools_list = tools

    def tools(self) -> list:
        return list(self.tools_list)

    def stop(self) -> None:
        for client in self.clients:
            try:
                client.__exit__(None, None, None)
            except Exception:
                pass
\end{lstlisting}

\texttt{GatewayToolPool} enters each client's context manager on initialisation and
registers a cleanup function with \texttt{atexit}. The Strands agent holds a reference
to this pool and calls \texttt{tools()} to get the current list. Because Strands holds
persistent MCP client connections, the pool must be kept alive for the duration of the
container's lifetime.

\paragraph{LangGraph gateway loader.}

\begin{lstlisting}[language=Python]
import os
from langchain_mcp_adapters.client import MultiServerMCPClient

async def load_gateway_tools() -> list:
    urls = _gateway_urls()
    if not urls:
        return []

    region = os.environ.get("GATEWAY_REGION", "us-west-2")
    client_config = {}
    for index, url in enumerate(urls):
        client_config[f"backend_{index}"] = {
            "transport": "http",
            "url": url,
            "auth": SigV4HTTPAuth(
                service=SERVICE_NAME, region=region
            ),
            "headers": {
                "Accept": "application/json, text/event-stream"
            },
        }

    client = MultiServerMCPClient(client_config)
    return await client.get_tools()
\end{lstlisting}

The LangGraph loader is a coroutine. It awaits \texttt{get\_tools()} and returns
LangChain-compatible tool objects that can be passed directly to
\texttt{create\_react\_agent}. Unlike the Strands pool,
\texttt{MultiServerMCPClient} is stateless by default: each tool invocation creates a
fresh MCP session, executes the tool, and closes the session. This means the LangGraph
implementation does not require session affinity at the container level, but it incurs a
small per-call connection overhead.

\section{Step 4: Encode operating procedures in SKILL.md}

Tools tell the agent \emph{how} to call a capability. A skill tells the agent \emph{when}
and \emph{why} to use capabilities in a particular order. The assistant has one skill: a
Markdown document that defines the standard operating procedure for the tasks it is
responsible for.

The skill document starts with YAML frontmatter that names the skill and declares which
gateway tools it is allowed to use:

\begin{lstlisting}
---
name: data-processing-assistant
description: >
  Standard operating procedure for data retrieval,
  job execution, and result review.
allowed-tools:
  - backend___fetch_data
  - backend___run_job
  - backend___check_status
  - backend___get_results
  - backend___create_override
---

## Processing Flow

1. Confirm the input identifier with the user.
2. Fetch the relevant data with fetch_data.
3. Trigger the processing job with run_job.
4. Poll completion status with check_status.
5. Retrieve and summarise results with get_results.
\end{lstlisting}

The skill registry reads these files and returns metadata for the system prompt and full
content for the agent's operating context:

\begin{lstlisting}[language=Python]
from pathlib import Path
from typing import Any

def skill_files(base_dir: Path | None = None) -> list[Path]:
    root = base_dir or Path(__file__).parent
    return sorted((root / "skills").glob("*/SKILL.md"))

def parse_skill_frontmatter(
    skill_file: Path,
    base_dir: Path | None = None,
) -> dict[str, Any]:
    text = skill_file.read_text(encoding="utf-8")
    root = base_dir or Path(__file__).parent
    metadata: dict[str, Any] = {
        "name": skill_file.parent.name,
        "description": "",
        "allowedTools": [],
        "path": str(skill_file.relative_to(root)),
    }
    if not text.startswith("---"):
        return metadata
    end = text.find("\n---", 3)
    if end == -1:
        return metadata
    current_key = ""
    for raw_line in text[3:end].splitlines():
        line = raw_line.rstrip()
        if not line.strip():
            continue
        if (line.startswith("  - ")
                and current_key == "allowed-tools"):
            metadata["allowedTools"].append(line[4:].strip())
            continue
        if ":" in line and not line.startswith(" "):
            key, value = line.split(":", 1)
            current_key = key.strip()
            if current_key == "name":
                metadata["name"] = value.strip()
            elif current_key == "description":
                metadata["description"] = value.strip()
            elif current_key == "allowed-tools":
                metadata["allowedTools"] = []
    return metadata

def loaded_skills(base_dir: Path | None = None) -> list[dict]:
    root = base_dir or Path(__file__).parent
    return [
        parse_skill_frontmatter(f, root)
        for f in skill_files(root)
    ]
\end{lstlisting}

The frontmatter parser keeps the logic minimal: it reads YAML keys without a YAML library
so the agent package stays lean. Strands consumes the skills directory through a plugin.
LangGraph reads the full document content and embeds it in the system prompt.

\section{Step 5: Build the Strands agent}

Strands takes a declarative approach. The developer provides a model, a session manager,
plugins, tools, and a system prompt. The framework owns the agent loop: it decides when to
call the model, when to invoke tools, and how to stream results back. The agent's code
describes \emph{what} the agent is, not \emph{how} its loop runs.

\begin{lstlisting}[language=Python]
import json, os
from pathlib import Path
from strands import Agent, AgentSkills
from bedrock_agentcore.memory.integrations.strands\
    .session_manager import AgentCoreMemorySessionManager

def _model_id() -> str:
    return os.environ.get(
        "MODEL_ID", "us.amazon.nova-2-lite-v1:0"
    )

_agent: Agent | None = None
_gateway_pool: GatewayToolPool | None = None

def _gateway_tools() -> list:
    global _gateway_pool
    if _gateway_pool is None:
        _gateway_pool = GatewayToolPool()
    tools = _gateway_pool.tools()
    if not tools:
        raise RuntimeError("No gateway tools loaded")
    return tools

def _skills_plugin():
    return AgentSkills(
        skills=str(Path(__file__).parent / "skills")
    )

def _memory_session_manager(
    runtime_session_id: str,
) -> AgentCoreMemorySessionManager:
    actor_id, session_id = _parse_runtime_session_id(
        runtime_session_id
    )
    # Replace with your AgentCore Memory ID and region
    return AgentCoreMemorySessionManager(
        agentcore_memory_config={
            "memory_id": os.environ["MEMORY_ID"],
            "actor_id": actor_id,
            "session_id": session_id,
        },
        region_name=os.environ.get(
            "MEMORY_REGION", "us-west-2"
        ),
    )

def _agent_instance(runtime_session_id: str) -> Agent:
    global _agent
    if _agent is None:
        gateway_tools = _gateway_tools()
        _agent = Agent(
            model=_model_id(),
            session_manager=_memory_session_manager(
                runtime_session_id
            ),
            plugins=[_skills_plugin()],
            tools=[
                get_current_time,
                echo,
                *gateway_tools,
            ],
            system_prompt=(
                "You are a helpful assistant. "
                "Use the available gateway tools to complete "
                "tasks. Follow the loaded skills for standard "
                "operating procedures.\n\n"
                f"Loaded skills: "
                f"{json.dumps(loaded_skills())}\n\n"
            ),
            callback_handler=None,
        )
    return _agent
\end{lstlisting}

Three things are worth observing here.

First, memory is attached as a \emph{session manager}. The Strands \texttt{Agent} receives
an \texttt{AgentCoreMemorySessionManager} and from that point treats memory as a framework
service. It retrieves relevant history, maintains conversational continuity, and writes new
interactions to memory as part of the normal agent loop. The application code does not
manage memory explicitly.

Second, skills arrive through the \texttt{AgentSkills} plugin. The plugin scans the
\texttt{skills/} directory, parses each \texttt{SKILL.md}, and makes its procedures
available to the agent without the developer having to embed content manually in a prompt.

Third, the container is pinned to a single session. If the same container receives a
request for a different session identifier, the agent raises an error. This is a deliberate
trade-off: because Strands holds stateful MCP client connections through the
\texttt{GatewayToolPool}, reusing a container across sessions would corrupt session state.
AgentCore Runtime handles this constraint by routing requests to the correct container.

\section{Step 6: Build the LangGraph agent}

LangGraph makes the agent loop explicit. The developer builds a graph with nodes, edges,
and hooks. The prebuilt \texttt{create\_react\_agent} supplies a default ReAct
loop---model call, tool call, observation, repeat---but every part of the loop can be
extended, replaced, or instrumented.

\paragraph{Memory components.}

LangGraph splits memory into two objects. The checkpointer stores graph state per thread
so conversations can resume after interruption. The store is a semantic store for
long-term memory that supports vector search over past interactions.

\begin{lstlisting}[language=Python]
from langgraph_checkpoint_aws import (
    AgentCoreMemorySaver,
    AgentCoreMemoryStore,
)

def _memory_components():
    memory_id = os.environ["MEMORY_ID"]
    region = os.environ.get("MEMORY_REGION", "us-west-2")
    checkpointer = AgentCoreMemorySaver(
        memory_id, region_name=region,
    )
    store = AgentCoreMemoryStore(
        memory_id=memory_id, region_name=region,
    )
    return checkpointer, store
\end{lstlisting}

\paragraph{The pre-model hook.}

The pre-model hook runs before every model call. It is the point where the graph reads
long-term memory and injects it into the conversation. LangGraph exposes this hook
explicitly because the developer may want to log what was retrieved, filter it, or change
the injection strategy for different graph branches. Strands handles equivalent retrieval
inside the framework.

\begin{lstlisting}[language=Python]
import uuid, json
from langchain_core.messages import HumanMessage
from langchain_core.runnables import RunnableConfig
from langgraph.store.base import BaseStore

def _pre_model_hook(
    state: dict,
    config: RunnableConfig,
    *,
    store: BaseStore,
) -> dict:
    configurable = config.get("configurable", {})
    actor_id = str(
        configurable.get("actor_id", "default_user")
    )
    thread_id = str(
        configurable.get("thread_id", "default")
    )
    messages = list(state.get("messages", []))

    latest_human = next(
        (
            m for m in reversed(messages)
            if isinstance(m, HumanMessage)
            or getattr(m, "type", None) == "human"
        ),
        None,
    )

    if latest_human is not None:
        store.put(
            (actor_id, thread_id),
            str(uuid.uuid4()),
            {"message": latest_human},
        )
        memories = store.search(
            (actor_id, thread_id),
            query=str(
                getattr(latest_human, "content", "")
            ),
            limit=5,
        )
        if memories:
            memory_lines = [
                json.dumps(
                    getattr(m, "value", m), default=str
                )
                for m in memories
            ]
            messages.insert(
                0,
                HumanMessage(
                    content=(
                        "Relevant memory context:\n"
                        + "\n".join(memory_lines)
                    )
                ),
            )

    return {"llm_input_messages": messages}
\end{lstlisting}

The hook saves the latest human message to the store and immediately searches the same
namespace for relevant prior context. If any is found, it is prepended to the message list
before the model sees the conversation. This is explicit: the developer controls exactly
when memory is written, what is searched, and how results are formatted before injection.

\paragraph{Assembling the graph.}

With tools, memory, and hooks in place, the graph is assembled in a single call:

\begin{lstlisting}[language=Python]
from langchain.chat_models import init_chat_model
from langgraph.prebuilt import create_react_agent

_graph = None
_gateway_tools_cache = None

def _build_system_prompt(gateway_tools: list) -> str:
    tool_names = [
        getattr(t, "name", str(t)) for t in gateway_tools
    ]
    return (
        "You are a helpful assistant. "
        "Use the available gateway tools to complete tasks. "
        "Follow the loaded skills for standard operating "
        "procedures.\n\n"
        f"Available gateway tools: {tool_names}\n\n"
        f"Loaded skills: {json.dumps(loaded_skills())}\n\n"
    )

async def _graph_instance():
    global _graph, _gateway_tools_cache
    if _graph is None:
        _gateway_tools_cache = await load_gateway_tools()
        if not _gateway_tools_cache:
            raise RuntimeError("No gateway tools loaded")

        checkpointer, store = _memory_components()
        region = os.environ.get("AWS_REGION", "us-west-2")
        model = init_chat_model(
            _model_id(),
            model_provider="bedrock_converse",
            region_name=region,
        )
        tools = [
            get_current_time,
            echo,
            *_gateway_tools_cache,
        ]
        _graph = create_react_agent(
            model=model,
            tools=tools,
            checkpointer=checkpointer,
            store=store,
            prompt=_build_system_prompt(_gateway_tools_cache),
            pre_model_hook=_pre_model_hook,
        )
    return _graph
\end{lstlisting}

The \texttt{create\_react\_agent} call builds a prebuilt ReAct graph: the model decides
which tools to call, each tool is executed, and the observation is fed back to the model
until it produces a final answer. Skills in LangGraph are embedded in the system prompt:
the skill registry reads each \texttt{SKILL.md} document and the
\texttt{\_build\_system\_prompt} function formats them into a single prompt block. There
is no plugin; the developer controls the skill injection point directly.

\section{Step 7: The entrypoint and session identity}

Both implementations wrap the agent in a \texttt{BedrockAgentCoreApp} and expose an
\texttt{@app.entrypoint} function. The entrypoint is the interface between the AgentCore
Runtime and the agent's domain logic. Both use the same convention: the
\texttt{runtimeSessionId} encodes both actor and session using a triple-underscore
separator.

\begin{lstlisting}[language=Python]
def _parse_runtime_session_id(raw: str) -> tuple[str, str]:
    if "___" in raw:
        actor_id, session_id = raw.split("___", 1)
        return actor_id.strip(), session_id.strip()
    return "default_user", raw.strip()

def _runtime_session_id_from(payload: dict, context) -> str:
    user_id = str(payload.get("user_id") or "").strip()
    session_id = str(payload.get("session_id") or "").strip()
    if user_id and session_id:
        return f"{user_id}___{session_id}"
    context_session_id = getattr(context, "session_id", None)
    if context_session_id:
        return str(context_session_id)
    raise ValueError(
        "Provide user_id + session_id in the payload, "
        "or invoke through AgentCore with a runtimeSessionId"
    )
\end{lstlisting}

\paragraph{Strands entrypoint.}

The Strands entrypoint creates or retrieves the agent instance, then calls
\texttt{stream\_async} with the user prompt. The framework handles tool calling, memory
updates, and streaming events internally.

\begin{lstlisting}[language=Python]
from bedrock_agentcore import BedrockAgentCoreApp

app = BedrockAgentCoreApp()

@app.entrypoint
async def invoke(payload: dict, context=None):
    prompt = str(payload.get("prompt") or "").strip()
    if not prompt:
        raise ValueError("prompt is required")
    if prompt.lower() in {"health", "ping", "status"}:
        yield {"message": {"content": [{"text": "ok"}]}}
        return
    runtime_session_id = _runtime_session_id_from(
        payload, context
    )
    agent = _agent_instance(runtime_session_id)
    async for event in agent.stream_async(prompt):
        yield event
\end{lstlisting}

\paragraph{LangGraph entrypoint.}

The LangGraph entrypoint awaits the graph, invokes it with the prompt and session
configuration, and yields a single response event extracted from the final messages list.

\begin{lstlisting}[language=Python]
from langchain_core.messages import AIMessage

app = BedrockAgentCoreApp()

def _last_ai_text(messages: list) -> str:
    for message in reversed(messages):
        if (
            isinstance(message, AIMessage)
            or getattr(message, "type", None) == "ai"
        ):
            content = getattr(message, "content", "")
            if isinstance(content, str):
                return content
            if isinstance(content, list):
                return "".join(
                    item.get("text", "")
                    if isinstance(item, dict)
                    else str(item)
                    for item in content
                )
    return ""

@app.entrypoint
async def invoke(payload: dict, context=None):
    prompt = str(payload.get("prompt") or "").strip()
    if not prompt:
        raise ValueError("prompt is required")
    if prompt.lower() in {"health", "ping", "status"}:
        yield {"message": {"content": [{"text": "ok"}]}}
        return
    runtime_session_id = _runtime_session_id_from(
        payload, context
    )
    actor_id, session_id = _parse_runtime_session_id(
        runtime_session_id
    )
    graph = await _graph_instance()
    result = await graph.ainvoke(
        {"messages": [("human", prompt)]},
        config={
            "configurable": {
                "actor_id": actor_id,
                "thread_id": session_id,
            }
        },
    )
    text = _last_ai_text(list(result.get("messages", [])))
    yield {
        "message": {
            "role": "assistant",
            "content": [{"text": text}],
        }
    }
\end{lstlisting}

The Strands version yields events as they arrive from the framework. The LangGraph version
awaits the full graph result and yields one final response. This reflects a deeper
difference: Strands is streaming-first by design, while LangGraph's \texttt{ainvoke}
returns a completed state.

\section{Step 8: The full picture}

With every component defined, the two agents have the same external shape. Both accept a
prompt over the AgentCore Runtime boundary. Both authenticate to the same MCP gateway
using SigV4. Both load the same skill document. Both store and retrieve memory using the
same AgentCore Memory resource. The difference is in what the developer writes.

With Strands, the developer writes an agent \emph{declaration}. The framework owns the
loop, the tool dispatch, the memory integration, and the event streaming. The developer
chooses model, tools, plugins, and prompt. The control flow is implicit in those choices.

With LangGraph, the developer writes an agent \emph{graph}. The framework provides a
graph builder and a default ReAct loop, but the developer can add nodes, attach hooks, add
conditional edges, and instrument every transition. Memory integration is explicit: the
developer writes the hook that reads memory and decides how to inject it. Skills
integration is explicit: the developer reads the skill document and decides how to format
it in the prompt.

\section{Framework comparison}

\begin{longtable}{p{0.22\textwidth}p{0.36\textwidth}p{0.32\textwidth}}
\toprule
\textbf{Aspect} & \textbf{Strands} & \textbf{LangGraph} \\
\midrule
Agent loop
    & Owned by the framework; hidden from application code
    & Explicit graph; developer builds nodes and edges \\
\addlinespace
Gateway auth
    & \texttt{SigV4HTTPAuth} passed to \texttt{MCPClient}
      via \texttt{streamablehttp\_client}
    & \texttt{SigV4HTTPAuth} passed as \texttt{auth} field
      in \texttt{MultiServerMCPClient} config \\
\addlinespace
MCP connection
    & Persistent connections held in \texttt{GatewayToolPool};
      tools loaded once at startup
    & Stateless per tool call; each invocation opens a fresh
      MCP session \\
\addlinespace
Memory
    & \texttt{AgentCoreMemorySessionManager} as a framework
      service; retrieval handled internally
    & \texttt{AgentCoreMemorySaver} +
      \texttt{AgentCoreMemoryStore} + explicit pre-model hook \\
\addlinespace
Skills
    & \texttt{AgentSkills} plugin scans the \texttt{skills/}
      directory automatically
    & Developer reads \texttt{SKILL.md} files and embeds
      content in the system prompt \\
\addlinespace
Streaming
    & \texttt{stream\_async()} emits events as the loop runs
    & \texttt{ainvoke()} returns completed state; streaming
      requires \texttt{astream()} \\
\addlinespace
Session affinity
    & Required: container pinned to a single session because
      MCP connections are stateful
    & Not required: stateless MCP sessions allow the container
      to serve any session \\
\addlinespace
Control flow
    & Framework decides tool dispatch order
    & Developer can add conditional edges, interrupts,
      and custom nodes \\
\bottomrule
\end{longtable}

\section{When to choose each framework}

The practical question is where complexity lives in the system.

If complexity lives in the tools and operating procedures---the skill is detailed, the
gateway exposes many capabilities, the conversation requires careful tool
sequencing---Strands keeps the application code compact while the skill documents carry
the domain logic. The framework handles the loop; the developer handles the domain.

If complexity lives in the orchestration itself---the agent must coordinate deterministic
steps with model-driven steps, expose checkpoints for human approval, branch on
domain-specific conditions, or make memory injection visible for evaluation---LangGraph
provides explicit structure for representing and testing each part of the trajectory.

Both frameworks integrate with AgentCore Runtime, AgentCore Gateway, and AgentCore Memory.
The decision is about where control should live: in the tools and operating procedures, or
in the graph itself.

\section{Summary}

Both implementations accept a prompt over the AgentCore Runtime boundary, authenticate to
the MCP gateway with SigV4, load tools from the gateway, follow skill-encoded operating
procedures, and persist state through AgentCore Memory. The gateway authentication
pattern---signing each inbound request with the container's IAM execution role
credentials---is the standard approach for any AWS workload calling an IAM-protected
service endpoint, and applies equally to both frameworks. What differs is the agent loop.
Strands declares an agent and delegates loop control to the framework. LangGraph builds an
explicit graph with hooks, memory components, and a system prompt that includes skill
content directly. The choice between them is a question of where the developer wants to
own control: in the tools and operating procedures, or in the graph itself.